The ballot initiative is the most successful mechanism for redistricting reform in the United States. We examine three landmark state initiative campaigns---California's Propositions 11 (2008) and 20 (2010), Colorado's Amendments Y and Z (2018), and Michigan's Proposal 2 (2018)---to identify the common design elements that predict success, the political barriers that defeated or weakened alternative proposals, and the lessons applicable to future initiatives incorporating algorithmic requirements. We find that successful initiatives share five design properties: independent commission structure, explicit partisan-balance requirements for commissioners, transparent public engagement mandates, enumerated criteria with priority ordering, and robust constitutional entrenchment that limits legislative repeal. None of the three states included an algorithmic requirement in its initial initiative, but all three created infrastructure that could readily accommodate one. We develop model ballot initiative language for California that would add an algorithmic redistricting requirement---specifically, recursive bisection using the \texttt{bisect} platform---to the existing California Citizens Redistricting Commission (CCRC) framework established by Propositions 11 and 20. This model language preserves the CCRC's human deliberation while providing an algorithmic baseline that constrains manipulation and enables verification.
